\begin{afterword}
\summaryhead

The post argues that one can understate confidence as well as overstate it, and singles out one case: displaying neutrality or suspended judgment in order to look mature, impartial or above the fray. The examples are the author's parents, who answered theological questions by saying they had grown out of asking them; a school principal who says it does not matter who started a fight; and Great Powers who tell smaller parties to stop fighting at once. The author calls this ``pretending to be Wise.''

Quoting Paulo Freire, the post argues that washing one's hands of a conflict sides with the powerful: refusing to ask who started a fight helps aggressors. The author attributes the pose partly to the convenience of the neutral party and partly to status. Suspending judgment is sometimes rational, but ``neutrality is a definite judgment'' about the balance of evidence, and it can be wrong. A call for compromise on abortion is likewise a definite position. One may decline to spend limited resources on divisive issues, but should not treat declining as a mark of wisdom.

\responsehead

The post's central point is correct and useful. To say that the evidence is balanced is to make a claim about the evidence, which can be checked and can be wrong, and a call for compromise on abortion is a position like any other. The intelligent-design example shows why this matters: a show of balance can misstate where the evidence lies.

Most of the rest is about motives, and the motives are asserted rather than shown. The definition itself is a motive: neutrality displayed ``in order to signal'' wisdom. The post gives no way to tell that from sincere neutrality, the same gap as in ``Professing and Cheering'' and ``Belief as Attire.'' The post attributes the neutrality of principals and Great Powers to convenience and status. The evidence is that a quick truce suits them, and that a principal who took sides would look like a participant. That a policy suits someone does not show that they adopt it for that reason, and no case is given in which it is shown.

The examples are weaker than they look. The principal's sentence has a second half, ``it only matters who ends it,'' which gives the children a reason not to hit back; the post discusses only the first half. The Great Powers get ``I believe'' and one sarcastic line about Israel and Hamas, posted a month after the Gaza War, with no demand for a truce named. The parents' case was told the day before, in ``Against Maturity,'' where the author explained the same reply as a story the parents told themselves after giving up on such questions, not as a display for others.

The post also contradicts the author's own earlier advice. ``Politics is the Mind-Killer'' said that a point about rationality should avoid contemporary politics ``if you can possibly avoid it,'' and should use Louis XVI if the point is political. This post uses abortion and the recent Gaza fighting, though its playground example needs neither, and then repeats the advice to practice on less charged topics.

In short: a correct point, that neutrality is itself a judgment, surrounded by motives assigned to principals, parents and Great Powers without evidence, and by political examples the author's own earlier advice would have avoided.
\end{afterword}
